MHC-I molecule-peptide binding prediction is a key step in tumor neoantigen identification, immunotherapy, and immunopeptidomics. With the candidate peptide space in the human proteome continuing to expand, deploying deep learning models for large-scale screening often runs into practical bottlenecks: high-dimensional peptide representations demand large caches, storage costs are high, I/O pressure is heavy, and online queries become slow. In this study, we take FennOmix.MHC as our base model and focus on model lightweighting and system-level deployment optimization for massive peptide sets.

We first reproduced the model and established an evaluation pipeline. Then we compared schemes with different output dimensions by looking at ranking performance, how well the embedding space structure is preserved, and the resulting cache footprint. A 64-dimensional embedding reduced the peptide cache from 132 GB to 19 GB while largely maintaining predictive accuracy-a reasonable trade-off between effectiveness and storage cost. Next, we examined two dimensionality reduction architectures, Linear and MLP. Taking into account standard prediction metrics and how they affect a downstream reduced-database search workflow, we adopted the simpler Linear projection as the default, as it is cheaper to deploy. Analyses of individual attention heads and feed-forward network units further revealed that the model contains redundancy; moderate pruning left predictive performance essentially unchanged.

On the system side, we designed an INT8-based quantization strategy for cache compression. Together with a custom binary cache format, hash indexing, and a combination of full loading and lazy loading, this allowed us to build a query pipeline that connects offline database construction, index lookup, cache retrieval, real-time encoding fallback, and online matching. Experimental results show that these steps further ease the burden of cache storage and loading, improving the deployment feasibility for screening very large peptide libraries. Taken together, we optimized FennOmix.MHC along several directions-representation compression, structural simplification, cache quantization, and index-based retrieval-providing a practical path for the engineering deployment of deep learning models in immune peptide binding prediction.